\originalchapter{红移、光度距离与宇宙学观测}
本章对应8.902第11讲及第12讲开头的共动体积。需要区分两个问题：宇宙中相隔多远，以及我们从收到的光推断出哪一种距离。膨胀空间里这两个问题不能用一个欧氏距离同时回答。沿用式\eqref{eq:astro-flrw}的无量纲$a_0=1$、长度坐标$\chi$和曲率$K$。

\originalsection{红移如何从两束相邻光波得到}
定义光谱红移$1+z=\lambda_{\rm obs}/\lambda_{\rm em}=\nu_{\rm em}/\nu_{\rm obs}$。先考虑固定共动坐标的光源与观测者，忽略其本动速度和局部引力势。径向光线沿零测地线传播，$\dd s^2=0$给出$c\dd t=a\dd\chi$。相邻两个波峰跨过同一段共动距离，故
\[
 \int_{t_e}^{t_o}\frac{c\dd t}{a(t)}
 =\int_{t_e+\delta t_e}^{t_o+\delta t_o}\frac{c\dd t}{a(t)}.
\]
当一周期远短于膨胀时间时，端点变化给$\delta t_o/a_o=\delta t_e/a_e$，于是
\begin{equation}\label{eq:astro-redshift}
 1+z=\frac{a_o}{a_e},\qquad a_e=\frac1{1+z}\quad(a_o=1).
\end{equation}
同一关系既降低光子能量，又拉长源的变化时间：$E_o=E_e/(1+z)$、$\delta t_o=(1+z)\delta t_e$。源的超新星光变因此发生时间伸长。光子沿零测地线传播并不表示空间切片中的路径总是普通意义的最短路。

狭义相对论的径向退行速度$\beta c$给$1+z_{\rm Dop}=\sqrt{(1+\beta)/(1-\beta)}$：在源的静止系两次发射相隔$\Delta\tau$，观测者系相隔$\gamma\Delta\tau$，第二个波峰又需多走$\beta c\gamma\Delta\tau$，总接收间隔为$\gamma(1+\beta)\Delta\tau$。这解释近距离$z\simeq\beta$。宇宙学红移则积分了光传播期间不断改变的几何；远距离不能由一个全局惯性系的$\beta$解释。存在本动与引力红移时，局部频率比与背景膨胀的频率比相乘，而不是把大红移简单相加。

\originalsection{径向共动距离与回溯时间}
从$\dot a=aH$和式\eqref{eq:astro-redshift}可得
\[
 \frac{\dd z}{\dd t}=-(1+z)H(z),\quad
 E^2(z)=\Omega_{r,0}(1+z)^4+\Omega_{m,0}(1+z)^3
 +\Omega_{K,0}(1+z)^2+\Omega_{\Lambda,0}.
\]
发射与接收相隔的\textbf{回溯时间}为
\[
 t_0-t_e=\frac1{H_0}\int_0^z\frac{\dd z'}{(1+z')E(z')},
\]
径向共动距离则是
\begin{equation}\label{eq:astro-comoving-distance}
 D_C=\chi=\frac c{H_0}\int_0^z\frac{\dd z'}{E(z')}.
\end{equation}
额外的$1+z$有明确来源：回溯时间积分的是物理时间，共动距离把每段光程除以当时的$a$。所以$c(t_0-t_e)$一般不是$D_C$。今天同一宇宙时间切片上的径向固有距离为$a_0\chi=\chi$，发射时则为$a_e\chi$；我们不能在同一时刻同时看见这条切片上的所有点。

共动距离由观测和模型推断，而不是禁止观测的量。其优势是理想共动物体的$\chi$不随宇宙膨胀改变。非零曲率时，径向距离$D_C$和横向共动距离
\[
 D_M=S_K(D_C)
\]
不同；下列角度和通量关系使用$D_M$。将$\Omega_K$与$K$的符号混淆会把正弦和双曲正弦用反。

\originalsection{用尺寸与通量分别定义距离}
设源在发射切片上的横向物理长度为$l$，小角尺寸为$\vartheta$。度规给$l=a_eS_K(\chi)\vartheta$，所以\textbf{角直径距离}定义及结果是
\[
 D_A=\frac l\vartheta=\frac{D_M}{1+z}.
\]
$l$必须指源当时的实际物理尺度，不能误用今日共动尺度。远处角直径距离可能随红移增加而减小；物体更远并不总是更小。

\textbf{光度距离}由源静止系的总光度$L$与观测总辐射通量$F$定义$F=L/(4\pi D_L^2)$。观测光束铺在今日面积$4\pi D_M^2$上；每个光子的能量和单位接收时间的光子数各减小一个$1+z$，故
\begin{equation}\label{eq:astro-distance-duality}
 F=\frac{L}{4\pi D_M^2(1+z)^2},\qquad
 D_L=(1+z)D_M=(1+z)^2D_A.
\end{equation}
后一个式子称距离对偶关系。这里假定光子数沿传播守恒、光走度规的零测地线，且用一致的源与光束面积定义；尘埃吸收需单独校正。引力透镜改变单源视通量和视尺寸，不能无视它给单个源距离估计造成的误差。

式\eqref{eq:astro-distance-duality}使用\emph{全频率积分}的通量。固定观测波段对应源不同的静止系波段，需星系章节讨论的谱形与$K$修正。例如以每单位频率表示，在相应频率$\nu_e=(1+z)\nu_o$处
\[
 F_{\nu_o}=\frac{(1+z)L_{\nu_e}}{4\pi D_L^2}.
\]
其积分$\dd\nu_e=(1+z)\dd\nu_o$恢复总通量式；写成每单位波长时因变量雅可比不同而有相反的额外因子。这是检验谱通量红移系数的简便办法。

\begin{example}
\textbf{Einstein--de Sitter宇宙的三个距离。}
取$\Omega_{m,0}=1$且其他参数为零，$E(z)=(1+z)^{3/2}$。直接积分得
\[
 \frac{H_0D_C}{c}=2\left[1-(1+z)^{-1/2}\right],\qquad
 D_A=\frac{D_C}{1+z},\qquad D_L=(1+z)D_C.
\]
当$z\ll1$三者都趋于$cz/H_0$。令$y=1+z$，$D_A\propto y^{-1}-y^{-3/2}$，导数为零要求$\sqrt y=3/2$，即$z=5/4$。这给角直径距离的最大点，而$D_L$仍随$z$单调增加。
\end{example}
\begin{center}
\begin{tikzpicture}
\begin{axis}[width=.76\textwidth,height=5.3cm,xmin=0,xmax=4,ymin=0,ymax=6,
 xlabel={$z$},ylabel={$H_0D/c$},legend pos=north west,grid=major]
\addplot[structurecolor,thick,domain=0:4,samples=100]{2*(1-1/sqrt(1+x))/(1+x)};
\addlegendentry{$D_A$}
\addplot[black,dashed,domain=0:4,samples=100]{2*(1-1/sqrt(1+x))};\addlegendentry{$D_C$}
\addplot[structurecolor,densely dotted,thick,domain=0:4,samples=100]{2*(1+x)*(1-1/sqrt(1+x))};\addlegendentry{$D_L$}
\end{axis}
\end{tikzpicture}
\end{center}
图由上述解析式自绘，展示的是物质单组分玩具模型，不是将课程中的第三方观测图重新嵌入。

\originalsection{天体计数需要的共动体积}
在红移宽度$\dd z$与立体角$\dd\Omega$内，共动横截面积是$D_M^2\dd\Omega$，径向厚度由式\eqref{eq:astro-comoving-distance}给出$c\dd z/H(z)$，所以
\[
 \frac{\dd V_c}{\dd z\dd\Omega}=\frac{cD_M^2}{H(z)}
 =\frac{c(1+z)^2D_A^2}{H(z)}.
\]
发射时固有体积再乘$a_e^3$。原笔记体积式中的个别$d\chi/dz$和额外径向因子排印混乱，按面积乘径向厚度即可辨认。若共动数密度$n_c(z)$已知，理想计数为$\dd N=n_c\dd V_c$；实际还需通量选择函数、遮挡、红移误差和天体演化。源静止系发生率$\mathcal R_c$给观测事件率时，还要除以$1+z$，不能把对象计数与事件率混为一式。

\originalsection{参数如何变成可比较的数据预测}
标准烛光先给距离模数$\mu=5\log_{10}(D_L/10\ \mathrm{pc})$，标准尺先给$\vartheta=l/D_A$。将测得的$z$与模型预测的$D_L,D_A$对比，可以约束式\eqref{eq:astro-Ea}的参数。绝对光度标定与$H_0$有简并，演化和消光能模仿距离变化，故单类数据不能未经系统误差处理就给宇宙结论。

低红移时定义减速参数$q_0=-\ddot a_0/(a_0H_0^2)$，Taylor展开给
\[
 D_L=\frac c{H_0}\left[z+\frac{1-q_0}{2}z^2+O(z^3)\right].
\]
一阶只能测膨胀率，二阶才携带加速信息；曲率在这一展开中从更高阶进入。高红移应积分完整方程，不能把二阶展开当成适用于$z\sim1$的精确关系。这一逻辑把均匀背景与实际超新星观测连接起来。
